\documentclass[a4paper, 11pt]{article}
\usepackage{graphicx}
\usepackage{amsmath,array}
%\usepackage{tablefootnote}
\usepackage{epstopdf}
\author{Yuxi Chen}
\title{Algorithm}
\begin{document}
\maketitle
\newpage

\section{Unit conversion}

A physical quantity contains two parts: the value times the unit. For example, the weight of a table is $w = 50 kg = 50,000 g$, where 50 and 50,000 are the values, and kg and g are the units. If we define a new unit wg, which is defined as $1 wg = 2.5 kg$, then the table weight is $w = 20 wg$. In general, a physical quantity can be expressed as $u = \bar{u} u^*$, where $\bar{u}$ is the value/number and $u^*$ is the unit.  

\subsection{Mass normalization}
\label{sec:mass_norm}

With the CGS unit, the particle motion equation is
\begin{eqnarray}
\frac{dv}{dt} = \frac{q}{m}(E + \frac{v}{c} \times B) 
\label{eq:velocity_update}
\end{eqnarray}
In FLEKS, all quantities are normalized. The normalization should not introduce new constants. So, from eq.\ref{eq:velocity_update}, the following expression is obtaied
\begin{equation}
\frac{q_{cgs}^* E_{cgs}^* t_{cgs}^*}{m_{cgs}^* v_{cgs}^*} = 1.
\label{eq:E1}
\end{equation}
From Gauss's law $\nabla\cdot E = 4\pi \rho$, we obtain 
\begin{equation}
E_{cgs}^* = \rho_{cgs} ^*x_{cgs}^*,
\label{eq:E2}
\end{equation}
where $\rho_{cgs}^* = q_{cgs}^*/x_{cgs}^{*3}$ is the charge density instead of mass density. From eq.\ref{eq:E1} and eq.\ref{eq:E2}, we obtain 
\begin{equation}
q_{cgs}^*=v_{cgs}^* \sqrt{m_{cgs}^* x_{cgs}^*}.
\label{eq:q1}
\end{equation}

In the code, $\bar{q}/(\bar{m}\bar{c}) = 1$ is required for a proton, so the normalization units $q_{cgs}^*$ and $m_{cgs}^*$ must satisfy 
\begin{equation}
q_{cgs}^*/(m_{cgs}^*v_{cgs}^*) = q_{cgs,p}/(m_{cgs,p}c_{cgs}),     
\end{equation}
where $q_{cgs,p}$ and $m_{cgs,p}$ are proton charge and mass in CGS unit, respectively. With eq.\ref{eq:q1}, we obtain the normalization mass in CGS unit:
\begin{eqnarray}
m_{cgs}^* = x_{cgs}^*(\frac{c_{cgs}m_{cgs,p}}{q_{cgs,p}})^2.
\end{eqnarray}
We note that $m_{cgs}^*$ also consists of the \textit{value} and the unit \textit{g}, and the expression above calculates its \textit{value}. Its \textit{value} in SI unit is calculated inside BATSRUS::ModPIC.f90, where the proton charge and mass are known in SI units. We need to convert the expression above into SI unit. Again, $m_{cgs}^*$ and $m_{SI}^*$ are the \textit{values} of the normalization mass in CGS and SI unit, respectively. If the normalization mass is 1.5kg, then $m_{SI}^* = 1.5$ and $m_{SI}^* = 1500$. The conversion is $m_{cgs}^* = 1000m_{SI}^*$. For the charge unit, $1C = 3\times 10^9 esu$ assume the speed of light is $3 \times 10^8m/s$. So, we have the following expression:
\begin{eqnarray}
1000m_{SI}^* = 100*x_{SI}^*(\frac{100*c_{SI}*1000*m_{SI,p}}{3\times10^9 q_{SI,p}})^2 \\
m_{SI}^* = 10^7x_{SI}^*(\frac{m_{SI,p}}{q_{SI,p}})^2.
\end{eqnarray}

Once $m_{SI}^*$ is obtained, it is passed to FLEKS and converted to $m_{cgs}^*$. 

\subsection{Length and velocity normalization}

These two are free parameters. Since Maxwell's equations and the particle motion equation are solved with $c = 1$ in code units (Section~\ref{sec:light_speed}), the velocity unit $v_{cgs}^*$ is also the speed of light of the simulation, $c' = v_{cgs}^*$.

\subsection{Charge and current normalization}

$q_{cgs}^*=v_{cgs}^* \sqrt{m_{cgs}^* x_{cgs}^*}$, as it is shown in eq.\ref{eq:q1}. $j_{cgs}^* = q_{cgs}^*v_{cgs}^*/x_{cgs}^{*3}$

\subsection{B and E normalization}

B and E have the same unit in CGS. Substituting eq.\ref{eq:q1} into eq.\ref{eq:E1} to obtain $E_{cgs}^* = B_{cgs}^* = \sqrt{\frac{m_{cgs}^*}{x_{cgs}^*3}}v_{cgs}^*$ 

\subsection{The consequences of changing light speed}
\label{sec:light_speed}

\paragraph{Summary.} The speed of light of the simulation is the velocity unit, $c' = v_{cgs}^*$ (\texttt{uNorm}). Choosing $c' < c$ is equivalent to keeping all charges, masses and the vacuum permeability $\mu_0$ physical, and increasing the vacuum permittivity:
\begin{equation}
\epsilon_0' = \frac{1}{\mu_0 c'^2} = \epsilon_0 \left(\frac{c}{c'}\right)^2 .
\end{equation}
This is the only change. A quantity is modified if and only if it depends on $\epsilon_0$. The Lorentz force, gyromotion, inertial lengths, Alfv\'en speed and the MHD electric fields are kept. The plasma frequency, Debye length, charge separation, electric field energy and light waves are modified.

\paragraph{How $c'$ enters the code.}
\begin{itemize}
  \item Maxwell's equations (Section~\ref{sec:full_pic}) and the particle motion equation are solved with $c = 1$ in code units. A reduced light speed is therefore not a separate parameter: choosing a small velocity unit \emph{is} choosing a small $c'$. In a stand-alone run, $c'$ is \texttt{uNorm} of \texttt{\#NORMALIZATION}. In a coupled run, it is the velocity unit passed from the MHD code.
  \item The mass unit is computed with the \emph{physical} $c$ (Section~\ref{sec:mass_norm}). A proton with $\bar{q} = \bar{m} = 1$ therefore has the physical $q/(mc)$, i.e.\ the physical gyrofrequency. With $\bar{c} \equiv c/c'$, the true speed of light in code units, this is the condition $\bar{q}/(\bar{m}\bar{c}) = 1$.
  \item As a consequence, the code charge of a particle is $c'/c$ times its true charge, while its code mass is the true mass. For the magnetic force this is exactly compensated by $1/c'$ in place of $1/c$. For Gauss's law it is not: the same number of particles carries $c'/c$ times less charge in code units.
  \item The MHD quantities $\rho$, $\mathbf{u}$, $p$ and $\mathbf{B}$ are converted to code units without any reference to $c$. The conversions of $\mathbf{J}$ and $\mathbf{E}$ use $c'$ (\texttt{Si2NoJ} and \texttt{Si2NoE} in \texttt{FluidInterface.cpp}). In SI language this is the $\mu_0$-fixed, $\epsilon_0'$-increased system above.
\end{itemize}

\paragraph{What is changed.}

\begin{center}
\begin{tabular}{p{0.38\textwidth} p{0.22\textwidth} p{0.32\textwidth}}
\hline
Quantity & Change for $c \to c'$ & Reason \\
\hline
Electric field produced by a given charge separation & $\times\,(c'/c)^2$ & Gauss's law, $\nabla\cdot\mathbf{E} = \rho_q/\epsilon_0'$ \\
$\delta n/n$ needed to sustain a given $\mathbf{E}$ & $\times\,(c/c')^2$ & same \\
Plasma frequency $\omega_{ps}$ & $\times\, c'/c$ & $\omega_{ps}^2 = n_s q_s^2/(\epsilon_0' m_s)$, or $\omega_{ps} = c'/d_s$ \\
$\omega_{pe}/\Omega_{ce}$ & $= c'/v_{Ae}$ & $v_{Ae} = B/\sqrt{\mu_0 n m_e}$ \\
Debye length $\lambda_D$ & $\times\, c/c'$ & $\lambda_D/d_e = v_{th,e}/c'$ \\
Displacement current $\epsilon_0'\,\partial\mathbf{E}/\partial t$ & $\times\,(c/c')^2$ & Amp\`ere's law \\
Vacuum EM waves & speed $c'$ & \\
Electric field energy density $\epsilon_0' E^2/2$ & $\times\,(c/c')^2$ for the same $\mathbf{E}$ & \\
Relativistic effects & $\gamma$ computed with $c'$ & only for relativistic test particles \\
\hline
\end{tabular}
\end{center}

Langmuir waves, Debye shielding, electron holes, double layers and other electrostatic or charge-separation phenomena therefore have modified frequencies, scales or amplitudes.

\paragraph{What is kept.}

\begin{center}
\begin{tabular}{p{0.45\textwidth} p{0.47\textwidth}}
\hline
Quantity & Reason \\
\hline
Lorentz force $q(\mathbf{E} + \mathbf{v}\times\mathbf{B})$ for given SI fields; $\mathbf{E}\times\mathbf{B}$ drift & no $\epsilon_0$; the $\mathbf{E}$ conversion and the force law are both exact \\
Gyrofrequency $\Omega_{cs}$, gyroradius $r_{Ls}$ & $qB/m$ is physical \\
Alfv\'en, sound and thermal speeds & depend only on $B$, $\rho$, $p$ and $\mu_0$ \\
Inertial lengths $d_s = c/\omega_{ps} = \sqrt{m_s/(\mu_0 n_s q_s^2)}$ & $\epsilon_0$ cancels; equivalently $d_i = v_A/\Omega_{ci}$ \\
Faraday's law; $\mu_0\mathbf{J} = \nabla\times\mathbf{B}$ when the displacement current is small & no $\epsilon_0$ \\
Motional, Hall and $\nabla\cdot\mathbf{P}_e$ electric fields & set by the electron force balance (generalized Ohm's law); requires quasi-neutrality \\
Poynting flux $\mathbf{E}\times\mathbf{B}/\mu_0$ & no $\epsilon_0$ \\
\hline
\end{tabular}
\end{center}

A common misreading is that the electric force is $\bar{c}$ times weaker, because the pusher uses $d\bar{\mathbf{v}}/d\bar{t} = \bar{\mathbf{E}}$ instead of $\bar{c}\,\bar{\mathbf{E}}$ for a proton. That compares Gaussian-unit fields defined with two different speeds of light. In Gaussian units $E$ itself scales as $1/c$, and this compensates the smaller code charge exactly. For a given electric field in SI units, the particle acceleration is physical. The ratio between the electric force and the $\mathbf{v}\times\mathbf{B}$ force is not altered. What is weakened is the electric field that a given charge separation produces, which is why $\omega_{pe}$ is lowered.

\paragraph{Why quasi-neutral physics survives.} In a quasi-neutral plasma the electric field is whatever the electrons need for force balance. The plasma supplies it through a tiny charge separation, not through the vacuum response. From Gauss's law, a field $E \sim uB$ that varies on a scale $L$ needs
\begin{equation}
\frac{\delta n}{n} \sim \frac{u}{v_A}\,\frac{d_i}{L}\left(\frac{v_A}{c'}\right)^2 .
\end{equation}
Reducing $c'$ raises $\delta n/n$, but the electric field is unchanged as long as $\delta n/n \ll 1$. The constraint is tightest at electron scales ($L \sim d_e$, electron flows $\sim v_{Ae}$), where $\delta n/n \sim (v_{Ae}/c')^2$.

\paragraph{Why reduce $c$ in FLEKS?} The reason is \emph{not} the light-wave CFL or the $\omega_{pe}\Delta t$ limit of explicit PIC:
\begin{itemize}
  \item The semi-implicit ($\theta \ge 0.5$) energy-conserving field solve is stable for $c'\Delta t/\Delta x > 1$ and $\omega_{pe}\Delta t \gg 1$.
  \item The full-PIC time step is set only by the particle CFL, $\Delta t = \mathrm{cfl}\cdot\Delta x/u_{\max}$, where $u_{\max}$ is the maximum thermal speed (\texttt{Pic.cpp}). The speed of light does not enter.
  \item The light CFL $c'\Delta t_{\text{sub}}\sqrt{\sum_d \Delta x_d^{-2}} < 1$ applies only to the explicit field solver and its field subcycling (Section~\ref{sec:full_pic}).
\end{itemize}
The cost of a large $c'$ shows up in the implicit linear system. Both the $\delta^2\nabla\times\nabla\times$ term ($\delta = c'\theta\Delta t$) and the mass-matrix term (of order $(\omega_{pe}\theta\Delta t)^2$) grow as $c'^2$. The system becomes ill-conditioned, GMRES needs many more iterations, and the converged field is more sensitive to particle noise. A smaller $c'$ keeps this system well conditioned.

\paragraph{Choosing $c'$.} $c'$ should exceed every physical speed by a comfortable margin, typically a factor of 10 or more:
\begin{enumerate}
  \item $c' \gg u$, $v_A$, $v_{th,i}$.
  \item $c' \gg v_{Ae}$ (equivalently $\omega_{pe} \gg \Omega_{ce}$). This keeps quasi-neutrality at electron scales, which matters for Hall fields and the electron diffusion region.
  \item $c' \gg v_{th,e}$. This keeps $\lambda_D \ll d_e$.
  \item $c'$ well above the fastest whistler phase speed to be modeled, $\approx v_{Ae}/2$ at $k d_e \approx 1$. Otherwise the full-PIC R-mode deviates from the Hall-MHD whistler. In \texttt{tests/whistler}, $c'/v_A \approx 1$, so the full-PIC run measures the full cold-plasma R-mode rather than the Hall whistler.
\end{enumerate}
$v_{Ae} = v_A\sqrt{m_i/m_e}$ uses the \emph{simulated} mass ratio. A reduced mass ratio and a reduced $c'$ therefore have to be chosen together. For example, with $v_A = 100$ km/s and $m_i/m_e = 100$, $v_{Ae} = 1000$ km/s, and $c' = 3000$ km/s gives only $\omega_{pe}/\Omega_{ce} = 3$.

\paragraph{Interpreting the output.}
\begin{itemize}
  \item $\mathbf{E}$, $\mathbf{B}$, $\mathbf{J}$, the moments and the Poynting flux converted to SI are physical, provided the plasma stays quasi-neutral.
  \item The net charge density (or $\nabla\cdot\mathbf{E}$), the electric field energy, and electrostatic oscillations and structures follow the ``changed'' table above. They should not be compared with observations without rescaling.
  \item The full-PIC mover is non-relativistic, so $c'$ is not a speed limit for the plasma particles. Particles faster than $c'$ indicate that $c'$ is too small.
  \item Relativistic test particles (\texttt{\#TPRELATIVISTIC}) compute $\gamma$ with the code speed of light, i.e.\ $c'$, and stop with an error when a particle exceeds $c'$. Their relativistic corrections are physical only if $c'$ is the true speed of light.
  \item The model is appropriate for non-relativistic, low-frequency ($\omega \ll \omega_{pe}'$), quasi-neutral physics dominated by magnetic and inductive effects. It shares this regime with Darwin and magneto-inductive models, but unlike them it retains the (slowed) displacement current and light waves.
\end{itemize}

\section{Particle mover}

\subsection{Boris algorithm}

FLEKS uses Boris particle mover: 
\begin{eqnarray}
\frac{\mathbf{v}^{n+1} - \mathbf{v}^{n}}{\Delta t} = \frac{q}{m}(\mathbf{E}^{n+\theta} + \bar{\mathbf{v}}\times\mathbf{B}) \\
\frac{\mathbf{x}^{n+1/2} - \mathbf{x}^{n-1/2}}{\Delta t} = \mathbf{v}^n
\end{eqnarray}
where $\mathbf{v}$ is defined as $\bar{\mathbf{v}} = \frac{\mathbf{v}^{n+1} + \mathbf{v}^{n}}{2}$. The traditional method updates the velocity with the following steps.

\begin{enumerate}
    \item Acceleration 
\begin{eqnarray}
    \mathbf{v}^- = \mathbf{v}^{n} + \frac{q\Delta t}{2m}\mathbf{E}^{n+\theta}
\end{eqnarray}

    \item Rotation
\begin{eqnarray} 
    \mathbf{a} = \mathbf{v}^- + \mathbf{v}^- \times \mathbf{t} \\
    \mathbf{v}^+ = \mathbf{v}^- + \mathbf{a} \times \mathbf{s}
\end{eqnarray}
    where $\mathbf{t} = \frac{q\Delta t}{2m} \mathbf{B}$, and $\mathbf{s} = \frac{2\mathbf{t}}{1+t^2}$.

    \item  Acceleration
\begin{eqnarray}    
    \mathbf{v}^{n+1} = \mathbf{v}^+ + \frac{q\Delta t}{2m}\mathbf{E}^{n+\theta}
\end{eqnarray}
\end{enumerate}

FLEKS solves the same velocity equation, but the implementation is slightly different.
\begin{enumerate}
    \item Acceleration 
\begin{eqnarray}
    \mathbf{v}^- = \mathbf{v}^{n} + \frac{q\Delta t}{2m}\mathbf{E}^{n+\theta}
\end{eqnarray}

\item Calculate $\bar{\mathbf{v}}$. 

With the expressions of $\mathbf{v}^+$ and $\mathbf{v}^-$ above, it is easy (it is probably not straightforwad, but I do not know how to number equations in markdown) to obtain $\mathbf{v}^{n+1} - \mathbf{v}^{n} = \mathbf{v}^+ - \mathbf{v}^- + \frac{q\Delta t}{m} \mathbf{E}^{n+\theta}$. Substitute it into the velocity update equation to remove electric field from the expression $\frac{\mathbf{v}^{+} - \mathbf{v}^{-}}{\Delta t} = \frac{q}{m}\bar{\mathbf{v}}\times\mathbf{B}$. Since $\mathbf{v}^+ = 2\bar{\mathbf{v}} - \mathbf{v}^-$, we obtain the equation for $\bar{\mathbf{v}}$:
\begin{eqnarray}
  \bar{\mathbf{v}} = \mathbf{v}^- + \frac{q\Delta t}{2m} \bar{\mathbf{v}} \times \mathbf{B} = \mathbf{v}^- + \bar{\mathbf{v}} \times \mathbf{t}   
\end{eqnarray}

To solve this equation, we apply $\cdot \mathbf{t}$ to the expression of $\bar{\mathbf{v}}$ to obtain $\bar{\mathbf{v}} \cdot \mathbf{t} = \mathbf{v}^- \cdot \mathbf{t}$. Then we apply $\times \mathbf{t}$ to the expression of $\bar{\mathbf{v}}$ to obtain
\begin{eqnarray}
\bar{\mathbf{v}} \times \mathbf{t} = \mathbf{v}^- \times \mathbf{t} + \bar{\mathbf{v}} \times \mathbf{t} \times \mathbf{t} = \mathbf{v}^- \times \mathbf{t} + (\bar{\mathbf{v}} \cdot \mathbf{t})\mathbf{t} - t^2\bar{\mathbf{v}} = \mathbf{v}^- \times \mathbf{t} + (\mathbf{v}^- \cdot \mathbf{t})\mathbf{t} - t^2\bar{\mathbf{v}}
\end{eqnarray}
Substitute the expression of $\bar{\mathbf{v}} \times \mathbf{t}$ to the expression of $\bar{\mathbf{v}}$ to obtain 
\begin{eqnarray}
\bar{\mathbf{v}} = \frac{\mathbf{v}^- + \mathbf{v}^- \times \mathbf{t} + (\mathbf{v}^- \cdot \mathbf{t})\mathbf{t}}{1+t^2}
\end{eqnarray}

\item Then it is easy to obtain $\mathbf{v}^{n+1} = 2*\bar{\mathbf{v}} - \mathbf{v}^{n}$
\end{enumerate}

\subsection{Relativistic Boris algorithm}

With $\mathbf{u} = \gamma \mathbf{v}$, the velocity update equation is
\begin{eqnarray}
\frac{\mathbf{u}^{n+1} - \mathbf{u}^{n}}{\Delta t} = \frac{q}{m}(\mathbf{E}^{n+\theta} + \bar{\mathbf{u}}\times \frac{\mathbf{B}}{\gamma ^{n+1/2}})
\end{eqnarray}

\begin{enumerate}
    \item Convert $\mathbf{v}$ to $\mathbf{u}$
    \item Acceleration
\begin{eqnarray}
    \mathbf{u}^- = \mathbf{u}^{n} + \frac{q\Delta t}{2m}\mathbf{E}^{n+\theta}
\end{eqnarray}

\item Calculate $\bar{\mathbf{u}}$ with
\begin{eqnarray}
\bar{\mathbf{u}} = \frac{\mathbf{u}^- + \mathbf{u}^- \times \mathbf{t} + (\mathbf{u}^- \cdot \mathbf{t})\mathbf{t}}{1+t^2}
\end{eqnarray}
where $\mathbf{t} = \frac{q\Delta t}{2m} \mathbf{B}/\gamma^{n+1/2}$. From the definition of $\gamma$ and $\mathbf{u} = \gamma \mathbf{v}$, $\gamma ^{n+1/2} = \sqrt{1+(u^-/c)^2}$ is obtained. 

\item Then it is easy to obtain $\mathbf{u}^{n+1} = 2*\bar{\mathbf{u}} - \mathbf{u}^{n}$
\item  Convert $\mathbf{u}$ back to $\mathbf{v}$
\end{enumerate}

\subsection{Boris mover linearization}

The Boris velocity update can be decomposed into an electric acceleration, a magnetic rotation, and a second electric acceleration. Denoting $\mathbf{t} = \frac{q\Delta t}{2m}\mathbf{B}^n$, the average velocity $\bar{\mathbf{v}} = \frac{\mathbf{v}^{n+1} + \mathbf{v}^n}{2}$ is related to the effective electric field $\mathbf{E}^{n+\theta}$ by:
\begin{equation}
\bar{\mathbf{v}} = \hat{\mathbf{v}} + \boldsymbol{\alpha}_p \cdot \mathbf{E}^{n+\theta},
\label{eq:vbar_linearized}
\end{equation}
where $\hat{\mathbf{v}}$ is the tentative velocity depending exclusively on the known velocity $\mathbf{v}^n$ and magnetic field $\mathbf{B}^n$:
\begin{equation}
\hat{\mathbf{v}} = \frac{\mathbf{v}^n + \mathbf{v}^n \times \mathbf{t} + (\mathbf{v}^n \cdot \mathbf{t})\mathbf{t}}{1 + t^2},
\end{equation}
and $\boldsymbol{\alpha}_p$ is the single-particle susceptibility tensor:
\begin{equation}
\boldsymbol{\alpha}_p = \frac{q\Delta t}{2m(1 + t^2)} \left[ \mathbf{I} + \mathbf{t}\mathbf{t} + \mathbf{I} \times \mathbf{t} \right].
\end{equation}
Here, $(\mathbf{I} \times \mathbf{t}) \cdot \mathbf{E} = \mathbf{E} \times \mathbf{t}$. Equation~\eqref{eq:vbar_linearized} provides the exact linear closure that couples particle kinetics to the implicit electromagnetic field equations in the full PIC semi-implicit method.

\section{Full PIC semi-implicit solver}
\label{sec:full_pic}

In the full-PIC approach, both electrons and ions are treated as kinetic particle species. To bypass the explicit Courant--Friedrichs--Lewy (CFL) constraint associated with the speed of light and plasma frequency ($\omega_{pe}\Delta t \le 2$), FLEKS employs the semi-implicit particle-in-cell method.

\subsection{Maxwell--Vlasov semi-implicit discretization}

Maxwell's equations in normalized units ($c=1$) are discretized in time using a $\theta$-implicit parameter ($\theta \in [0.5, 1.0]$):
\begin{align}
\frac{\mathbf{B}^{n+1} - \mathbf{B}^n}{\Delta t} &= -\nabla \times \mathbf{E}^{n+\theta}, \label{eq:maxwell_faraday} \\
\frac{\mathbf{E}^{n+1} - \mathbf{E}^n}{\Delta t} &= \nabla \times \mathbf{B}^{n+\theta} - 4\pi \bar{\mathbf{J}}, \label{eq:maxwell_ampere}
\end{align}
where intermediate fields are evaluated as:
\begin{align}
\mathbf{E}^{n+\theta} &= \theta \mathbf{E}^{n+1} + (1-\theta)\mathbf{E}^n, \\
\mathbf{B}^{n+\theta} &= \theta \mathbf{B}^{n+1} + (1-\theta)\mathbf{B}^n = \mathbf{B}^n - \theta \Delta t \nabla \times \mathbf{E}^{n+\theta}.
\end{align}
Substituting $\mathbf{B}^{n+\theta}$ into Amp\`ere's law~\eqref{eq:maxwell_ampere} yields:
\begin{equation}
\mathbf{E}^{n+\theta} - \mathbf{E}^n = \theta \Delta t \left[ \nabla \times \left(\mathbf{B}^n - \theta \Delta t \nabla \times \mathbf{E}^{n+\theta}\right) - 4\pi \bar{\mathbf{J}} \right].
\label{eq:ampere_intermediate}
\end{equation}

\subsection{Susceptibility mass matrix and tentative current}

Using the linearized Boris mover relation~\eqref{eq:vbar_linearized}, the average current density $\bar{\mathbf{J}}(\mathbf{x})$ is expressed as:
\begin{equation}
\bar{\mathbf{J}}(\mathbf{x}) = \hat{\mathbf{J}}(\mathbf{x}) + \mathbf{M} \cdot \mathbf{E}^{n+\theta},
\end{equation}
where $\hat{\mathbf{J}}(\mathbf{x})$ is the tentative current density deposited onto grid nodes using shape functions $S(\mathbf{x} - \mathbf{x}_p)$:
\begin{equation}
\hat{\mathbf{J}}(\mathbf{x}) = \sum_p q_p \hat{\mathbf{v}}_p S(\mathbf{x} - \mathbf{x}_p),
\end{equation}
and $\mathbf{M}$ is the nodal susceptibility mass matrix:
\begin{equation}
\mathbf{M}(\mathbf{x}, \mathbf{x}') = \sum_p q_p S(\mathbf{x} - \mathbf{x}_p) \boldsymbol{\alpha}_p S(\mathbf{x}' - \mathbf{x}_p).
\end{equation}
The mass matrix $\mathbf{M}$ is accumulated at node locations with support across neighboring stencil points, supporting block-adaptive mesh refinement (AMR).

\subsection{Implicit electric field linear system}

Substituting $\bar{\mathbf{J}}$ into eq.~\eqref{eq:ampere_intermediate} and rearranging gives the implicit linear system for $\mathbf{E}^{n+\theta}$:
\begin{equation}
\left[ \mathbf{I} + \theta^2 \Delta t^2 \nabla \times (\nabla \times \cdot) + 4\pi \theta \Delta t \mathbf{M} \right] \mathbf{E}^{n+\theta} = \mathbf{R},
\label{eq:esolve_operator}
\end{equation}
where the right-hand side $\mathbf{R}$ is evaluated from the known state at $t^n$:
\begin{equation}
\mathbf{R} = \mathbf{E}^n + \theta \Delta t \left( \nabla \times \mathbf{B}^n - 4\pi \hat{\mathbf{J}} \right).
\end{equation}
Using the vector identity $\nabla \times (\nabla \times \mathbf{E}) = \nabla (\nabla \cdot \mathbf{E}) - \nabla^2 \mathbf{E}$, the linear operator is implemented matrix-free and solved iteratively using GMRES or Jacobian-Free Newton--Krylov.

Once $\mathbf{E}^{n+\theta}$ is obtained, the fields are advanced:
\begin{align}
\mathbf{E}^{n+1} &= \frac{1}{\theta}\mathbf{E}^{n+\theta} - \frac{1-\theta}{\theta}\mathbf{E}^n, \\
\mathbf{B}^{n+1} &= \mathbf{B}^n - \Delta t \nabla \times \mathbf{E}^{n+\theta}.
\end{align}
To maintain Gauss's law $\nabla \cdot \mathbf{E} = 4\pi \rho_c$, an optional div($\mathbf{E}$) cleaning step solves the Poisson equation $\nabla^2 \phi = \frac{1}{\theta}(\nabla \cdot \mathbf{E} - 4\pi \rho_c)$ and applies a localized particle position correction $\Delta \mathbf{x} \propto -\nabla \phi$ or electric field projection.

\section{Upwind schemes and comoving frame solving}

In simulations featuring high-speed plasma flows, shocks, and sharp boundary layers, standard central-difference discretizations can produce unphysical dispersion, numerical oscillations, and grid heating. FLEKS incorporates a comoving frame formulation and directional upwind artificial viscosity schemes for both the electric and magnetic fields. (Note: The schemes described in Section~\ref{sec:upwind_E} and Section~\ref{sec:upwind_B} apply to the full-PIC semi-implicit solver; the hybrid PIC solver employs the nodal Rusanov dissipation described in Section~\ref{sec:hybrid_upwind}.)

\subsection{Comoving frame field solving}

When plasmas possess large bulk drift velocities (such as in the supersonic solar wind), solving fields in the stationary laboratory frame Doppler-shifts high-frequency wave modes and requires resolving fast advection. FLEKS provides the option to solve the electromagnetic fields in the frame comoving with the plasma bulk flow:

\begin{enumerate}
  \item The background fluid velocity $\mathbf{U}_0$ is computed from the ion fluid moments deposited at grid nodes:
  \begin{equation}
  \mathbf{U}_0(\mathbf{x}) = \frac{\sum_s (\rho \mathbf{u})_s}{\sum_s \rho_s}.
  \end{equation}
  Optionally, $\mathbf{U}_0$ may be smoothed with a digital filter to suppress PIC shot noise.
  \item The motional electric field in the comoving frame is:
  \begin{equation}
  \mathbf{E}_0 = -\mathbf{U}_0 \times \mathbf{B}^n.
  \end{equation}
  \item The tentative particle velocities and tentative current $\hat{\mathbf{J}}$ are evaluated relative to the comoving frame: $\hat{\mathbf{v}}' = \hat{\mathbf{v}} - \mathbf{U}_0$.
  \item The motional electric field generates an effective source current through the susceptibility matrix. Its contribution $\mathbf{M} \cdot \mathbf{E}_0$ is added to the right-hand side of the implicit system~\eqref{eq:esolve_operator}:
  \begin{equation}
  \mathbf{R}_{\text{comov}} = \mathbf{R} + 4\pi \theta \Delta t \, \mathbf{M} \cdot \mathbf{E}_0.
  \end{equation}
\end{enumerate}

\subsection{Upwind scheme for electric field}
\label{sec:upwind_E}

A directional Lax--Friedrichs artificial viscosity is added to the semi-implicit electric field equation to suppress numerical oscillations near shocks and steep gradients.

The artificial diffusion flux at cell interface $i+1/2$ along spatial direction $d$ is:
\begin{equation}
F_{i+1/2}^{\text{art}} = \frac{1}{2} c_R u_R (E_{i+1} - E_i),
\end{equation}
where $u_R$ is the characteristic propagation velocity. Two choices are supported:
\begin{itemize}
  \item $u_R = c_{\max}$, the (normalized) speed of light, giving isotropic maximum-speed upwinding.
  \item $u_R = |(\mathbf{U}_0)_d|$, the local component of the background plasma velocity along the interface direction. This limits numerical dissipation to the actual advection velocity and is used in conjunction with the comoving frame solver.
\end{itemize}
The coefficient $c_R$ is governed by a generalized Minmod/Sweby flux limiter:
\begin{equation}
c_R = 1 - \phi(r_R), \qquad \phi(r) = \max\left(0,\, \min\left(\theta_E r,\; \frac{1+r}{2},\; \theta_E\right)\right), \quad r = \frac{E_i - E_{i-1}}{E_{i+1} - E_i},
\end{equation}
where $\theta_E \in [0, 2]$ controls the amount of limiting:
\begin{itemize}
  \item $\theta_E = 0$: $\phi(r) = 0$, $c_R = 1$ --- fully diffusive first-order Lax--Friedrichs upwind scheme.
  \item $\theta_E = 1$: reproduces the minmod limiter.
  \item $\theta_E = 2$: least diffusive TVD limiter.
\end{itemize}

In the implicit electric field operator~\eqref{eq:esolve_operator}, this upwind term enters the left-hand side:
\begin{equation}
\Delta \mathcal{A}(\mathbf{E})_i = -\frac{\theta \Delta t}{2\Delta x} \left[ c_R u_R (E_{i+1} - E_i) - c_L u_L (E_i - E_{i-1}) \right].
\end{equation}

When the limiter coefficients $c_{R/L}$ are evaluated from the lagged field $\mathbf{E}^n$, the operator remains strictly linear and GMRES is used. When $c_{R/L}$ depends on the current iterate $\mathbf{E}$, the operator is non-linear and a matrix-free Newton--Krylov solver is applied instead.

\subsection{Upwind scheme for magnetic field}
\label{sec:upwind_B}

A directional upwind diffusion correction is applied to the cell-centered magnetic field following the Faraday update. The correction applied to cell $i$ is:
\begin{equation}
\Delta B_i = \frac{\Delta t}{2\Delta x} \left[ c_R u_R (B_{i+1} - B_i) - c_L u_L (B_i - B_{i-1}) \right].
\end{equation}
Key properties of the upwind $\mathbf{B}$ scheme include:
\begin{itemize}
  \item \textbf{Advection velocity:} $u_R$ and $u_L$ are face-reconstructed from the background plasma velocity $\mathbf{U}_0$. A constant advection speed can alternatively be prescribed.
  \item \textbf{Slope limiter:} $c_{R/L} = 1 - \phi(r)$ uses the same Minmod/Sweby limiter family as the electric field scheme, with an independent parameter $\theta_B \in [0, 2]$.
  \item \textbf{Boundary protection:} At physical domain boundaries, diffusion is suppressed when the flow is directed outward, preventing spurious boundary reflections.
\end{itemize}

\section{Hybrid PIC solver}

In the hybrid PIC approach, ions are treated as kinetic particles and electrons are modeled as a massless, charge-neutralizing fluid. The electromagnetic fields are evolved through a generalized Ohm's law rather than through the full implicit Maxwell solver.

\subsection{The generalized Ohm's law}

The electric field is assembled on the node grid from the node-centered ion fluid moments and the cell-centered magnetic field:
\begin{equation}
\mathbf{E} = -\mathbf{U}_i \times \mathbf{B}
           + \eta\,\mathbf{J}
           + \frac{\mathbf{J} \times \mathbf{B}}{\rho_q}
           - \frac{\nabla P_e}{\rho_q}
           + \mathbf{E}_{\text{upwind}}
           - \frac{\eta_h}{4\pi}\,\nabla \times (\nabla^2 \mathbf{B}),
\label{eq:ohmslaw}
\end{equation}
where:
\begin{itemize}
  \item $-\mathbf{U}_i \times \mathbf{B}$ is the \emph{convection} term, always active.
  \item $\eta\,\mathbf{J}$ is the \emph{resistive} term, with $\eta$ the magnetic diffusivity.
  \item $(\mathbf{J}\times\mathbf{B})/\rho_q$ is the \emph{Hall} term. When the Hall term is absent, the Hall-whistler time-step restriction is removed.
  \item $-\nabla P_e/\rho_q$ is the \emph{electron pressure gradient} term. The electron pressure $P_e$ supports two closures:
  \begin{itemize}
    \item Isothermal ($\gamma=1$): $P_e = T_e \rho$.
    \item Adiabatic ($\gamma > 1$): $P_e = P_0 (\rho/\rho_0)^\gamma$, where $P_0 = \rho_0 T_e$.
  \end{itemize}
  A minimum density floor $\rho_{\min}$ is applied to the $1/\rho_q$ factors in the Hall and electron pressure terms to prevent unbounded contributions in near-vacuum cells.
  \item $\mathbf{E}_{\text{upwind}}$ is the \emph{upwind (Rusanov) dissipation} term (Section~\ref{sec:hybrid_upwind}), enabled by default (\texttt{\#HYBRIDUPWIND}). It adds directional wave dissipation to suppress grid-scale oscillations while strictly preserving discrete $\nabla \cdot \mathbf{B} = 0$.
  \item $-(\eta_h/4\pi)\,\nabla\times(\nabla^2\mathbf{B})$ is the \emph{hyper-resistivity} term. Hyper-resistivity damps grid-scale PIC shot noise without introducing excessive dissipation at large scales. The coefficient $\eta_h$ can be set to scale with the grid resolution:
  \begin{equation}
  \eta_h = 4\pi C_h \frac{\Delta x_{\min}^4}{\Delta t_{\text{sub}}},
  \end{equation}
  where $C_h$ is a dimensionless tuning parameter, typically $10^{-3}$--$10^{-2}$.
\end{itemize}
The current density is obtained from Amp\`ere's law (CGS), $\mathbf{J} = \nabla\times\mathbf{B}/(4\pi)$. The explicit resistive and hyper-resistive diffusion terms must satisfy parabolic CFL stability bounds:
\begin{equation}
\frac{\eta k_{\max}^2 \Delta t_{\text{sub}}}{4\pi} \le C_{\text{CFL}}, \qquad \frac{\eta_h k_{\max}^4 \Delta t_{\text{sub}}}{4\pi} \le C_{\text{CFL}},
\end{equation}
where $C_{\text{CFL}} = 2.785$ for RK4 and $2.513$ for SSPRK3. The upwind dissipation term satisfies the hyperbolic CFL criterion discussed in Section~\ref{sec:hybrid_upwind}.

\subsection{Collocated field layout}

The hybrid solver uses the same grid layout as the full-PIC solver: the electric field $\mathbf{E}$, the current density $\mathbf{J}$, and the ion moments $\rho$, $\rho\mathbf{u}$ are stored on the node grid, while the magnetic field $\mathbf{B}$ and the scalar electron pressure $P_e$ are stored at cell centers. The generalized Ohm's law is therefore assembled at the nodes from the node-centered moments and the cell-centered $\mathbf{B}$, and the Faraday update advances the cell-centered $\mathbf{B}$; the node-centered $\mathbf{B}$ that the Boris push samples is recovered by averaging the surrounding cell values. The cell-centered field is differentiated with the compact $2\,\Delta x$ central-difference stencil
\begin{equation}
(\nabla \times \mathbf{B})_{i,j,k} \propto \frac{\mathbf{B}_{i+1} - \mathbf{B}_{i-1}}{2\Delta x},
\end{equation}
whose discrete Fourier kernel vanishes identically at the grid Nyquist mode ($k\Delta x = \pi$). This eliminates the grid-scale odd-even decoupling instability that would otherwise force heavy Hall sub-cycling.

\subsection{Time integration and subcycling of $\mathbf{B}$}

The magnetic field is advanced via Faraday's law:
\begin{equation}
\frac{\partial \mathbf{B}}{\partial t} = -\nabla\times\mathbf{E},
\end{equation}
using an explicit Runge--Kutta integrator. Two choices are provided:
\begin{itemize}
  \item \textbf{RK4}: Classical fourth-order Runge--Kutta. Because RK4 has a wide stability region along the imaginary axis (where Hall-whistler oscillations reside), a single subcycle per coarse step is often sufficient.
  \item \textbf{SSPRK3}: Strong-stability-preserving third-order Runge--Kutta with time-centered intermediate stages.
\end{itemize}
Within one coarse particle time step $\Delta t$, $N_{\text{sub}}$ subcycles of step size $\Delta t_{\text{sub}} = \Delta t / N_{\text{sub}}$ are taken.

To achieve second-order temporal accuracy, the ion moments are linearly interpolated across each magnetic sub-step fraction $h_{\text{step}} \in [0, 1]$ between the previous deposit $X^{n-1/2}$ and the current deposit $X^{n+1/2}$:
\begin{equation}
X(h_{\text{step}}) = \left(\tfrac{1}{2}-h_{\text{step}}\right) X^{n-1/2} + \left(\tfrac{1}{2}+h_{\text{step}}\right) X^{n+1/2}.
\end{equation}

\subsection{Upwind (Rusanov) dissipation}
\label{sec:hybrid_upwind}

Central-difference schemes applied to the generalized Ohm's law and Faraday's law are susceptible to odd--even grid-scale decoupling ($2\Delta x$ modes), macroparticle shot noise, and numerical oscillations across steep gradients or shocks. While hyper-resistivity can suppress short-wavelength noise, its explicit parabolic stability limit $\Delta t_{\text{sub}} \propto \Delta x^4 / \eta_h$ requires fine-tuned coefficients and can force severe subcycling on refined grids.

To provide robust, scale-aware dissipation without parabolic time-step penalties, FLEKS implements a directional upwind (Rusanov / local Lax--Friedrichs) dissipation scheme (\texttt{\#HYBRIDUPWIND}), which is enabled by default in the hybrid solver.

\subsubsection{Nodal Rusanov formulation and divergence preservation}

In standard finite-volume magnetohydrodynamics, Rusanov dissipation diffuses cell-centered fluxes. If applied directly as an artificial diffusion on the cell-centered magnetic field (as in the full-PIC \texttt{\#UPWINDB} option), directional dissipation violates the solenoidal constraint $\nabla \cdot \mathbf{B} = 0$, requiring hyperbolic or elliptic divergence cleaning.

In the hybrid PIC solver, FLEKS circumvents this limitation by injecting the dissipation directly into the node-centered electric field in Ohm's law:
\begin{equation}
\mathbf{E}_{\text{upwind}} = \sum_{d=1}^D \frac{a_d}{2}\, \hat{\mathbf{e}}_d \times \left(\mathbf{B}^R_d - \mathbf{B}^L_d\right),
\label{eq:hybrid_upwind}
\end{equation}
where $\hat{\mathbf{e}}_d$ is the unit vector along spatial direction $d \in \{x, y, z\}$, $a_d$ is the directional characteristic signal speed, and $\mathbf{B}^L_d, \mathbf{B}^R_d$ are the reconstructed magnetic field states on the low and high sides of the node along direction $d$.

When Faraday's law $\partial \mathbf{B}/\partial t = -\nabla \times \mathbf{E}$ is evaluated via the discrete node-to-cell curl $\operatorname{curl}_{\text{node}\to\text{center}}$, the contribution of $\mathbf{E}_{\text{upwind}}$ along direction $d$ becomes:
\begin{equation}
-\nabla \times \left[ \frac{a_d}{2}\, \hat{\mathbf{e}}_d \times \Delta\mathbf{B}_d \right], \qquad \Delta\mathbf{B}_d \equiv \mathbf{B}^R_d - \mathbf{B}^L_d.
\end{equation}
For example, along the $x$-direction ($d=x$, $\hat{\mathbf{e}}_x$):
\begin{equation}
\hat{\mathbf{e}}_x \times \Delta\mathbf{B}_x = -\Delta B_{x,z}\,\hat{\mathbf{e}}_y + \Delta B_{x,y}\,\hat{\mathbf{e}}_z,
\end{equation}
and taking the curl yields:
\begin{align}
(-\nabla \times \mathbf{E}_{\text{upwind}, x})_x &= 0, \\
(-\nabla \times \mathbf{E}_{\text{upwind}, x})_y &= \partial_x E_{\text{upwind}, z} = \partial_x \left( \frac{a_x}{2}\,\Delta B_{x,y} \right), \\
(-\nabla \times \mathbf{E}_{\text{upwind}, x})_z &= -\partial_x E_{\text{upwind}, y} = \partial_x \left( \frac{a_x}{2}\,\Delta B_{x,z} \right).
\end{align}
In a one-dimensional discrete stencil, the node-to-cell curl advances cell $j$ for each transverse component as:
\begin{equation}
\frac{\partial B_j}{\partial t} = \dots + \frac{a}{2\Delta x} \left( J_{j+1/2} - J_{j-1/2} \right),
\end{equation}
where $J_{j+1/2} = \mathbf{B}^R_{j+1/2} - \mathbf{B}^L_{j+1/2}$ is the jump across face $j+1/2$. This formulation guarantees three crucial properties:
\begin{enumerate}
  \item \textbf{Normal field invariance:} The longitudinal magnetic field component normal to the face ($\mathbf{B} \cdot \hat{\mathbf{e}}_d$) satisfies $\hat{\mathbf{e}}_d \times (\Delta B_d\,\hat{\mathbf{e}}_d) = 0$; hence, normal magnetic flux across cell interfaces is never diffused.
  \item \textbf{Machine-precision divergence preservation:} Because $\mathbf{E}_{\text{upwind}}$ is added directly to the nodal electric field, the discrete identity
  \begin{equation}
  \operatorname{div}_{\text{center}}\left(\operatorname{curl}_{\text{node}\to\text{center}} \mathbf{E}\right) \equiv 0
  \end{equation}
  holds unconditionally on uniform grids and across AMR hierarchies (where the standard coarse--fine interface electric field synchronization applies to $\mathbf{E}_{\text{upwind}}$ without modification). The update preserves $\nabla \cdot \mathbf{B} = 0$ to machine precision without requiring any divergence cleaning.
  \item \textbf{Intrinsic field preservation:} Jumps $\mathbf{B}^R_d - \mathbf{B}^L_d$ are reconstructed exclusively from the evolved magnetic field $\mathbf{B}$; static intrinsic dipole or background fields $\mathbf{B}_0$ are excluded, avoiding spurious diffusion of planetary dipoles.
\end{enumerate}

\subsubsection{MUSCL reconstruction and slope limiters}

The cell-centered magnetic field $\mathbf{B}$ is reconstructed to interface states using a high-resolution MUSCL approach. Across the face between cells $i-1$ and $i$ along direction $d$:
\begin{equation}
\mathbf{B}^L_{i-1/2} = \mathbf{B}_{i-1} + \frac{1}{2}\mathbf{s}_{i-1}, \qquad \mathbf{B}^R_{i-1/2} = \mathbf{B}_i - \frac{1}{2}\mathbf{s}_i,
\end{equation}
where $\mathbf{s}_i$ is the limited slope vector in cell $i$. For each field component, denoting consecutive differences $\Delta_L = B_i - B_{i-1}$ and $\Delta_R = B_{i+1} - B_i$, FLEKS supports four slope limiters (\texttt{upwindLimiter}):
\begin{itemize}
  \item \texttt{upwind1} (first-order): $\mathbf{s}_i = 0$, giving an unattenuated first-order jump $\mathbf{B}^R - \mathbf{B}^L = B_i - B_{i-1}$.
  \item \texttt{minmod}:
  \begin{equation}
  s = \operatorname{minmod}(\Delta_L, \Delta_R) = \begin{cases}
  \min(\Delta_L, \Delta_R), & \Delta_L > 0 \text{ and } \Delta_R > 0, \\
  \max(\Delta_L, \Delta_R), & \Delta_L < 0 \text{ and } \Delta_R < 0, \\
  0, & \Delta_L \Delta_R \le 0.
  \end{cases}
  \end{equation}
  \item \texttt{vanleer}:
  \begin{equation}
  s = \begin{cases}
  \dfrac{2\Delta_L \Delta_R}{\Delta_L + \Delta_R}, & \Delta_L \Delta_R > 0, \\
  0, & \Delta_L \Delta_R \le 0.
  \end{cases}
  \end{equation}
  \item \texttt{mc} (monotonized central, default):
  \begin{equation}
  s = \begin{cases}
  \operatorname{sgn}(\Delta_L)\min\left(2|\Delta_L|,\, 2|\Delta_R|,\, \frac{1}{2}|\Delta_L + \Delta_R|\right), & \Delta_L \Delta_R > 0, \\
  0, & \Delta_L \Delta_R \le 0.
  \end{cases}
  \end{equation}
\end{itemize}

\paragraph{Transverse face averaging.}
In 2D and 3D, nodes are surrounded by multiple cell faces perpendicular to direction $d$. To eliminate duplicate reconstructions and maintain identical stencil symmetry with $\operatorname{curl}_{\text{center}\to\text{node}}$, each face jump is computed once on the face grid and stored. The jump at node $(i,j,k)$ is then formed by averaging the surrounding face jumps along the transverse directions (over 2 faces in 2D, 4 faces in 3D):
\begin{equation}
\Delta\mathbf{B}_{d, \text{node}} = \frac{1}{2^{D-1}} \sum_{\text{transverse}} \Delta\mathbf{B}_{d, \text{face}}.
\end{equation}

\subsubsection{Characteristic signal speed}

The directional propagation speed $a_d$ governing the Rusanov dissipation is:
\begin{equation}
a_d = |u_{e,d}| + c_{\text{fast}} + \alpha_w\,\frac{|\mathbf{B}|}{4\rho_{\text{eff}}\Delta x_d},
\label{eq:hybrid_upwind_speed}
\end{equation}
where:
\begin{itemize}
  \item $\mathbf{u}_e = \mathbf{u}_i - \mathbf{J}/\rho_{\text{eff}}$ is the electron fluid velocity, accounting for Hall drift (reducing to ion bulk velocity $\mathbf{u}_i$ when the Hall term is disabled). Here $\rho_{\text{eff}} = \max(\rho, \rho_{\min})$ avoids singularities in depleted regions.
  \item $c_{\text{fast}} = \sqrt{v_A^2 + c_{s,e}^2} = \sqrt{\frac{|\mathbf{B}|^2}{4\pi\rho_{\text{eff}}} + \gamma_e T_e}$ is the fast magnetosonic speed based on the local magnetic field and electron temperature ($c_{s,e}^2 = \gamma_e T_e$).
  \item The final term represents the whistler wave phase speed at the grid cutoff wavenumber $k_{\max, d} = \pi / \Delta x_d$:
  \begin{equation}
  v_{\phi, w} = \frac{\omega_w}{k_{\max, d}} = \frac{k_{\max, d}|\mathbf{B}|}{4\pi\rho_{\text{eff}}} = \frac{\pi|\mathbf{B}|}{4\pi\rho_{\text{eff}}\Delta x_d} = \frac{|\mathbf{B}|}{4\rho_{\text{eff}}\Delta x_d},
  \end{equation}
  scaled by the dimensionless coefficient $\alpha_w = \texttt{upwindCoefWhistler}$ (default $0.0$).
\end{itemize}
Setting $\alpha_w = 0.0$ upwinds with the fluid advection and fast magnetosonic speed; $\alpha_w = 1.0$ yields the full Hall-MHD Rusanov speed. Because the whistler phase speed scales as $O(1/\Delta x)$, non-zero $\alpha_w$ introduces noticeable damping on resolved whistler modes ($7\%$ damping per radian at $\alpha_w = 1.0$ versus $0.2\%$ at $\alpha_w = 0.0$ for a mode resolved by 16 cells). The default $\alpha_w = 0.0$ provides stable dissipation of grid-scale noise across all standard hybrid tests without attenuating physical whistler waves.

For code verification and testing, the characteristic speed can be overridden by a user-specified constant $a_d = \texttt{upwindFixedSpeed} > 0$.

\subsubsection{Grid-scale smoothness sensor}

Away from local extrema, second-order limited schemes (such as MC) damp Fourier modes as $O(\Delta x^3 k^4)$. Although significantly less diffusive than first-order upwinding ($O(\Delta x k^2)$), this damping may accumulate over long simulation durations on well-resolved physical waves.

To restrict dissipation strictly to the grid scale without widening the stencil beyond the two ghost layers already present in the Faraday advance, FLEKS introduces a grid-scale smoothness sensor. Each component of the reconstructed face jump is modulated by a factor $\phi \in [\phi_{\text{floor}}, 1]$:
\begin{equation}
\widetilde{\Delta\mathbf{B}} = \phi\,\left(\mathbf{B}^R - \mathbf{B}^L\right), \qquad \phi = \min\left(1,\, \max\left(\phi_{\text{floor}},\, s\right)\right),
\end{equation}
where $\phi_{\text{floor}} = \texttt{upwindSensorFloor}$ (default $0.1$), and the sensor $s$ measures the normalized third spatial difference across the four adjacent cells $(i-2, i-1, i, i+1)$ spanning the face:
\begin{equation}
s = \frac{|B_{i+1} - 3B_i + 3B_{i-1} - B_{i-2}|}{|B_i - B_{i-1}| + |B_{i-1} - B_{i-2}| + |B_{i+1} - B_i|}.
\end{equation}
For a smooth resolved wave $B \propto e^{ikx}$, Taylor expansion gives:
\begin{equation}
s \approx \frac{(k\Delta x)^2}{4} \ll 1,
\end{equation}
causing $\phi$ to clamp to the floor $\phi_{\text{floor}} = 0.1$, which suppresses numerical damping on resolved waves by an order of magnitude (lowering damping on a 16-cell whistler from $0.2\%$ to $0.02\%$ per radian). At grid-scale Nyquist noise ($k\Delta x = \pi$, where $B_i \propto (-1)^i$) and near shock jumps, the third difference is $O(1)$, so $s \sim 4/3 > 1$, resulting in $\phi = 1$ and delivering full Rusanov dissipation. Setting $\phi_{\text{floor}} = 1.0$ disables the sensor.

\subsubsection{Stability and CFL condition}

The first-order upwind term damps the Nyquist mode ($k_d\Delta x_d = \pi$) at the real eigenvalue rate:
\begin{equation}
\gamma_d = \frac{2 a_d}{\Delta x_d}.
\end{equation}
The total dissipation rate across all active grid dimensions is:
\begin{equation}
\Omega_{\text{upwind}} = \sum_{d=1}^D \frac{2 a_d}{\Delta x_d}.
\end{equation}
Explicit Runge--Kutta stability along the negative real axis requires:
\begin{equation}
\Delta t_{\text{sub}}\,\Omega_{\text{upwind}} \le C_{\text{CFL}},
\end{equation}
where $C_{\text{CFL}} = 2.785$ for RK4 and $2.513$ for SSPRK3.

FLEKS evaluates $\Omega_{\text{upwind}}$ during time-step and subcycle planning (\texttt{Pic::calc\_hybrid\_dt\_and\_subcycle}). Applying a safety factor ($c_{\text{safety}} = 0.85$), the subcycle time step is constrained by:
\begin{equation}
\Delta t_{\text{sub}} \le \frac{c_{\text{safety}} C_{\text{CFL}}}{\Omega_{\text{upwind}}}.
\end{equation}
If this limit is exceeded, FLEKS issues a CFL warning and increases the subcycle count $N_{\text{sub}}$ accordingly. Because $\Omega_{\text{upwind}}$ scales as $O(1/\Delta x)$ rather than the parabolic $O(1/\Delta x^4)$ of hyper-resistivity, the upwind dissipation permits significantly larger time steps on fine AMR meshes while robustly maintaining stability.

\subsection{Algorithmic cycle and noise mitigation}

The complete hybrid PIC time step proceeds through the following stages:
\begin{enumerate}
  \item \textbf{Particle Boris push:} Ions are advanced $\mathbf{v}^n \to \mathbf{v}^{n+1}$ using integer-step fields $\mathbf{E}^n, \mathbf{B}^n$; positions are updated as $\mathbf{x}^{n+1/2} = \mathbf{x}^{n-1/2} + \Delta t\,\mathbf{v}^n$.
  \item \textbf{Kinetic source and boundary updates:} Particle boundary conditions, chemistry, charge exchange, and physical inflow flux injection are evaluated.
  \item \textbf{Moment deposition:} Ion charge $\rho_q^{n+1/2}$ and mass momentum $(\rho\mathbf{u})_i^{n+1/2}$ are deposited on the node grid. A digital-filter smoothing can be applied to reduce particle noise.
  \item \textbf{Faraday advance:} Subcycled explicit RK4 or SSPRK3 advances the cell-centered $\mathbf{B}^n \to \mathbf{B}^{n+1}$ using the interpolated ion moments and stage electric fields (evaluating Ohm's law with the upwind dissipation $\mathbf{E}_{\text{upwind}}$ at intermediate stages); the node-centered $\mathbf{B}^{n+1}$ sampled by the pusher is obtained by averaging the surrounding cell values.
  \item \textbf{Electric field assembly:} The node-centered $\mathbf{E}^{n+1}$ is computed from eq.~\eqref{eq:ohmslaw} (including $\mathbf{E}_{\text{upwind}}$) and stored for the subsequent Boris push.
\end{enumerate}

\subsection{Divergence-free magnetic field evolution with AMR}

\subsubsection{The divergence problem on AMR hierarchies}
On a single uniform grid, evaluating Faraday's law $\partial \mathbf{B}/\partial t = -\nabla\times\mathbf{E}$ via the discrete node-to-cell curl guarantees that the discrete divergence vanishes identically to machine precision:
\begin{equation}
\nabla \cdot (\nabla \times \mathbf{E}) \equiv 0.
\end{equation}
In multi-level block-structured Adaptive Mesh Refinement (AMR), however, standard approaches introduce severe divergence errors:
\begin{enumerate}
  \item \textbf{Interface flux mismatch:} When coarse and fine levels evaluate $\mathbf{E}$ independently, electric fields along coarse-fine boundaries are multi-valued, violating flux cancellation across shared cell faces and creating numerical magnetic monopoles.
  \item \textbf{Ghost cell interpolation:} Interpolating cell-centered $\mathbf{B}$ from coarse to fine ghost cells via standard polynomial/conservative interpolation does not satisfy the discrete divergence constraint, injecting divergence errors at every sub-step.
  \item \textbf{Covered coarse cell drift:} In block-structured AMR, coarse cells covered by fine patches cannot simply be replaced by the volume average of fine cells ($\mathbf{B}_{\text{coarse}} = \operatorname{avg}(\mathbf{B}_{\text{fine}})$), as direct restriction violates coarse $\nabla \cdot \mathbf{B} = 0$ along the boundary between covered and uncovered cells.
  Conversely, advancing covered coarse cells with their own coarse Faraday update leads to rapid divergence from the fine solution: the discrete operator $\operatorname{curl}_{\text{node}\to\text{center}}$ is non-telescoping across refinement levels ($2\Delta x$ coarse stencil vs. $\Delta x$ fine stencil). Even when coincident coarse nodes are given the exact fine electric field (injection), for any nodal $\mathbf{E}$
  \begin{equation}
  \operatorname{curl}_{\text{coarse}}(\mathbf{E}) \neq \operatorname{avg}(\operatorname{curl}_{\text{fine}}(\mathbf{E})),
  \end{equation}
  differing at $O(1)$ relative error throughout covered regions. Left uncorrected, covered coarse cells drift linearly away from the fine solution.
\end{enumerate}

\subsubsection{Interface synchronization and Faraday-evolved ghost layer}
FLEKS resolves interface divergence through a multi-pass synchronization:
\begin{enumerate}
  \item \textbf{Interface $\mathbf{E}$ synchronization:} Prior to evaluating $\nabla \times \mathbf{E}$, fine nodes on the coarse-fine interface are linearly interpolated along coarse edges from coarse nodal $\mathbf{E}$ (\texttt{sync\_emf\_at\_interface}), so both levels advance the cells sharing an interface face with the same $\mathbf{E}$. Fine ghost nodes are then filled from the coarse $\mathbf{E}$ (\texttt{fill\_ghost\_emf\_from\_coarse}). The covered coarse nodes keep the coarse $\mathbf{E}$: injecting the coincident fine values there does not make the covered coarse $\mathbf{B}$ follow the fine solution (see above), and it copies the grid-scale part of the fine $\mathbf{E}$ onto the coarse level, where it drives a weakly damped $2\Delta x_{\text{coarse}}$ mode along the interface.
  \item \textbf{Faraday-evolved first fine ghost layer:} Rather than interpolating all fine ghost cells from the coarse level, the first fine ghost layer is evolved directly via Faraday's law using the synchronized nodal $\mathbf{E}$ (\texttt{faraday\_ngrow > 0}). Only outer ghost layers ($\ge 2$) are interpolated from the coarse level.
  \item \textbf{True stage evaluation:} Intermediate RK stages evaluate Ohm's law directly at stage states $\mathbf{B}_{\text{stage}}$ with the synchronized $\mathbf{E}$ and Faraday-evolved ghost layer, eliminating ad-hoc lagged states to preserve high-order temporal accuracy.
\end{enumerate}

\subsubsection{Local divergence-free covered-cell relaxation}
To eliminate the drift of covered coarse cells without violating coarse divergence, FLEKS applies a \emph{local divergence-free relaxation} (\texttt{relax\_covered\_B\_to\_fine}):
\begin{equation}
\mathbf{B}_{\text{coarse}} \leftarrow \mathbf{B}_{\text{coarse}} + \nabla \times \mathbf{A},
\label{eq:relax_update}
\end{equation}
where $\mathbf{A}$ is a nodal potential defined strictly on interior nodes of the covered region (nodes whose surrounding cells are all refined). Boundary and uncovered nodes satisfy $\mathbf{A} \equiv 0$.

\paragraph{Divergence preservation:} Because $\operatorname{div}_{\text{node}\to\text{center}} \circ \operatorname{avg}_{\text{center}\to\text{node}} \circ \operatorname{curl}_{\text{node}\to\text{center}} \equiv 0$ holds identically in both 2D and 3D, and $\mathbf{A}$ vanishes outside the fully covered patch, the correction $\nabla \times \mathbf{A}$ has support strictly within covered cells and satisfies:
\begin{equation}
\nabla \cdot (\nabla \times \mathbf{A}) \equiv 0.
\end{equation}
Uncovered coarse cells and interface fluxes are completely untouched, preserving coarse-level $\nabla \cdot \mathbf{B} = 0$ to machine precision regardless of solver convergence.

\paragraph{Optimization and CGLS solve:} The potential $\mathbf{A}$ is determined by minimizing the least-squares mismatch between the coarse field and the fine cell average:
\begin{equation}
\min_{\mathbf{A}} \|\mathbf{r}\|^2 = \min_{\mathbf{A}} \|\operatorname{avg}(\mathbf{B}_{\text{fine}}) - \mathbf{B}_{\text{coarse}} - \nabla \times \mathbf{A}\|^2,
\end{equation}
yielding the normal equations:
\begin{equation}
(\nabla \times)^T (\nabla \times) \mathbf{A} = (\nabla \times)^T [\operatorname{avg}(\mathbf{B}_{\text{fine}}) - \mathbf{B}_{\text{coarse}}].
\end{equation}
This system is solved using the Conjugate Gradient for Least Squares (CGLS) method:
\begin{itemize}
  \item In 2D and pseudo-2D ($N_z = 1$), $B_z$ does not enter $\nabla \cdot \mathbf{B}$; it is restricted directly via cell averaging, while $(B_x, B_y)$ are relaxed using the scalar potential $A_z$.
  \item In true 3D, all three components are relaxed using the vector potential $\mathbf{A} = (A_x, A_y, A_z)$. Gauge transformations $\mathbf{A} \to \mathbf{A} + \nabla \phi$ lie in the null space of $\nabla \times$; CGLS initialized with $\mathbf{A}_0 = 0$ naturally operates orthogonal to the null space, eliminating the need for explicit gauge fixing.
\end{itemize}

\paragraph{Advantages and limitations:}
\begin{itemize}
  \item \emph{Advantages:} Strict machine-precision divergence preservation ($\sim 10^{-15}$); strictly local computation (no global elliptic Poisson solves); unconditional divergence safety under early iteration termination; and clean transmission of plasma waves across AMR interfaces.
  \item \emph{Limitations:} Nodes on the perimeter of the covered region are constrained to $\mathbf{A} = 0$, leaving a localized residual drift in the outermost covered layer; divergence-carrying numerical noise on the fine level cannot be represented by $\nabla \times \mathbf{A}$; and the Faraday-evolved first ghost layer can drift over long simulation times in strong-gradient regimes (monitored by the \texttt{dGhost} diagnostic).
\end{itemize}

\section{Explicit full PIC solver}

In the explicit full PIC solver (\texttt{\#FIELDSOLVER explicit}), both ions and electrons are treated as fully kinetic particles, and Maxwell's equations are evolved explicitly in time without requiring large-scale matrix inversions or Krylov solvers.

\subsection{Maxwell's equations and discrete dual staggering}

The electromagnetic fields satisfy Maxwell's equations:
\begin{align}
\frac{\partial \mathbf{B}}{\partial t} &= -c\,\nabla \times \mathbf{E}, \label{eq:faraday_explicit} \\
\frac{\partial \mathbf{E}}{\partial t} &= c\,\nabla \times \mathbf{B} - 4\pi \mathbf{J}, \label{eq:ampere_explicit}
\end{align}
with the divergence constraints $\nabla \cdot \mathbf{B} = 0$ and $\nabla \cdot \mathbf{E} = 4\pi \rho_c$.

To preserve identical multi-level MultiFab definitions across all FLEKS solvers, the explicit full PIC solver adopts a discrete dual-staggered arrangement:
\begin{itemize}
  \item The electric field $\mathbf{E}$ and current density $\mathbf{J}$ reside at \emph{grid nodes} (vertices).
  \item The magnetic field $\mathbf{B}$ resides at \emph{cell centers}.
\end{itemize}

The discrete spatial curl operators map between node-centered and cell-centered fields:
\begin{itemize}
  \item $\operatorname{curl}_{\text{node}\to\text{center}}(\mathbf{E})$ maps the vertex-centered electric field to cell centers via face circulation integrals:
  \begin{equation}
  (\nabla \times \mathbf{E})_{i+1/2, j+1/2, k+1/2} = \mathcal{C}_{\text{node}\to\text{center}}(\mathbf{E}).
  \end{equation}
  \item $\operatorname{curl}_{\text{center}\to\text{node}}(\mathbf{B})$ maps the cell-centered magnetic field to vertices via dual-cell circulations:
  \begin{equation}
  (\nabla \times \mathbf{B})_{i, j, k} = \mathcal{C}_{\text{center}\to\text{node}}(\mathbf{B}).
  \end{equation}
\end{itemize}
Because the discrete divergence of $\operatorname{curl}_{\text{node}\to\text{center}}$ identically vanishes to machine precision ($\operatorname{div}_{\text{center}}(\mathcal{C}_{\text{node}\to\text{center}}(\mathbf{E})) \equiv 0$), the discrete magnetic divergence constraint
\begin{equation}
\nabla \cdot \mathbf{B} = 0
\end{equation}
is strictly conserved over all time steps and across AMR coarse-fine interfaces without needing divergence cleaning potentials.

\subsection{Field subcycling and CFL condition}

Because the speed-of-light wave propagation imposes the explicit Courant--Friedrichs--Lewy (CFL) stability criterion:
\begin{equation}
c\,\Delta t_{\text{sub}} \sqrt{\sum_{d=1}^D \frac{1}{\Delta x_d^2}} < 1,
\end{equation}
the Maxwell equations can be subcycled within one particle coarse time step $\Delta t$. For $N_{\text{sub}}$ subcycles (\texttt{\#FIELDSUBCYCLE}), the subcycle time step is $\Delta t_{\text{sub}} = \Delta t / N_{\text{sub}}$. During each subcycle $m = 1, \dots, N_{\text{sub}}$, the field updates advance $\mathbf{B}$ and $\mathbf{E}$ while holding or interpolating the macroparticle current $\mathbf{J}$.

\subsection{Charge-conserving current}

The current is deposited as the particles move (Esirkepov's method, adapted to the FLEKS layout), so that the continuity equation, and with it Gauss's law, holds exactly on the grid. Gauss's law is evaluated at cell centers,
\begin{equation}
\operatorname{div}_{\text{node}\to\text{center}}(\mathbf{E}) = 4\pi \bar\rho_c, \qquad
\bar\rho_c = A_x A_y A_z\, \rho_{\text{node}},
\end{equation}
where $\rho_{\text{node}}$ is the node (CIC) charge and $A_d$ averages the two nodes of a cell along $d$, so a particle contributes to the cell centers with a trapezoidal shape three cells wide. A particle moving from $\mathbf{x}^0$ to $\mathbf{x}^1$ (at most one cell per dimension) with node shapes $S^0$, $S^1$ and $\Delta S = S^1 - S^0$ adds, along each dimension $d$ with transverse dimensions $e, f$, the current through the faces between the nodes
\begin{equation}
W_d = \Delta S_d \left( S^0_e S^0_f + \tfrac12 \Delta S_e S^0_f + \tfrac12 S^0_e \Delta S_f + \tfrac13 \Delta S_e \Delta S_f \right), \qquad
F_{d,i+1/2} = -\frac{q\,\Delta x_d}{\Delta t\, V} \sum_{i' \le i} W_{d,i'},
\end{equation}
and the node current is the average of the two adjacent faces, $J_{d,i} = (F_{d,i-1/2} + F_{d,i+1/2})/2$. Then
\begin{equation}
\bar\rho_c^{\,n+1} - \bar\rho_c^{\,n} = -\Delta t\,\operatorname{div}_{\text{node}\to\text{center}}(\mathbf{J})
\end{equation}
holds to round-off. Because $\operatorname{div}_{\text{node}\to\text{center}}(\operatorname{curl}_{\text{center}\to\text{node}}\mathbf{B}) \equiv 0$, the Amp\`ere update keeps $\operatorname{div}(\mathbf{E}) - 4\pi\bar\rho_c$ at its initial value, for any number of field subcycles. The cell-centered CIC charge cannot satisfy such a relation: its checkerboard modes lie in the null space of the node gradient. In 2D the out-of-plane component is $q\,v_z\,\bar S/V$ with the time-averaged shape $\bar S$.

The particles gather $\mathbf{E}$ with the transpose of this deposit: along its own direction $d$, $E_d$ is the average of the two nodes of the cell that holds the particle; across $d$ (and for $E_z$ in 2D) it is linear. With the plain CIC gather instead, the mismatch between the work done on the particles and the field energy heats a thermal plasma (the Langmuir test: $+0.8\%$ in $4\,$s against $+0.1\%$ with the matched gather). $\mathbf{B}$ is gathered with the CIC shape.

At a reflecting particle wall, the current of a particle is deposited along its unreflected path, and the part that falls on the ghost nodes beyond the wall is folded onto the mirror nodes with mirror parity (odd for the normal component, even for the tangential ones); the wall node is cleared for the normal component and doubled for the tangential ones, as the moments are. The folded current is that of the reflected path (image method), so the relation above holds in the wall cells as well, with the wall-node charge doubled as in the moments.

The div(E) correction of the semi-implicit solver is therefore off by default (\texttt{\#DIVE} turns it on). The conservation is exact for periodic domains and reflecting walls; particle resampling (\texttt{\#RESAMPLING}), boundary injection and removal, current smoothing (\texttt{\#SMOOTHJ}), embedded bodies and AMR coarse--fine interfaces change the charge or the current outside the deposit and leave a residual there. The div(E) correction targets the cell CIC charge rather than $\bar\rho_c$, so it does not repair that residual consistently.

\subsection{Time integrators}

Two time-advancement algorithms are provided via \texttt{\#FIELDINTEGRATOR}:
\begin{enumerate}
  \item \textbf{Boris--Maxwell leapfrog} (\texttt{leapfrog}, default):
  Particle positions and the fields $\mathbf{E}$, $\mathbf{B}$ are stored at integer steps $t^n$ and particle velocities at half-integer steps $t^{n+1/2}$. One step proceeds as:
  \begin{enumerate}
    \item Push: kick $\mathbf{v}^{n-1/2} \to \mathbf{v}^{n+1/2}$ with the Boris pusher and $(\mathbf{E}^n, \mathbf{B}^n)$ at $\mathbf{x}^n$, drift $\mathbf{x}^{n+1} = \mathbf{x}^n + \Delta t\,\mathbf{v}^{n+1/2}$, and deposit the charge-conserving current $\mathbf{J}^{n+1/2}$ of the displacement, all in one pass over the particles.
    \item Fields, within each subcycle $m$ (Strang splitting with $\mathbf{J}^{n+1/2}$ held fixed):
    \begin{align}
    \mathbf{B}^{m+1/2} &= \mathbf{B}^{m} - \tfrac{1}{2} c\,\Delta t_{\text{sub}}\,\operatorname{curl}_{\text{node}\to\text{center}}(\mathbf{E}^m), \\
    \mathbf{E}^{m+1} &= \mathbf{E}^m + c\,\Delta t_{\text{sub}}\,\operatorname{curl}_{\text{center}\to\text{node}}(\mathbf{B}^{m+1/2}) - 4\pi\,\Delta t_{\text{sub}}\,\mathbf{J}^{n+1/2}, \\
    \mathbf{B}^{m+1} &= \mathbf{B}^{m+1/2} - \tfrac{1}{2} c\,\Delta t_{\text{sub}}\,\operatorname{curl}_{\text{node}\to\text{center}}(\mathbf{E}^{m+1}).
    \end{align}
  \end{enumerate}
  For a cold plasma oscillation this scheme is neutrally stable for $\omega_{pe}\Delta t < 2$. A current centered at $t^n$ instead (for example $q\mathbf{v}^{n-1/2}$ deposited at $\mathbf{x}^n$) grows at $\gamma \approx \omega_{pe}^2 \Delta t / 4$.

  \item \textbf{Predictor-Corrector / Runge--Kutta 2} (\texttt{pc} or \texttt{rk2}):
  Velocities live at integer steps and particles drift with the average of the old and new velocities.
  \begin{itemize}
    \item \emph{Predictor}: a trial push with $(\mathbf{E}^n, \mathbf{B}^n)$ deposits the current $\mathbf{J}^*$ of the trial displacement without changing the particles; the fields advance to $(\mathbf{E}^*, \mathbf{B}^*)$ at $t^{n+1}$ with $\mathbf{J}^*$.
    \item \emph{Corrector}: the particles are pushed from $t^n$ with the time-centered fields $(\mathbf{E}^n + \mathbf{E}^*)/2$, $(\mathbf{B}^n + \mathbf{B}^*)/2$, depositing $\mathbf{J}^{n+1/2}$; the fields advance from $(\mathbf{E}^n, \mathbf{B}^n)$ to $t^{n+1}$ with it.
  \end{itemize}
  Since the predictor leaves the particles in place, no particle copy is made.
\end{enumerate}

\subsection{Memory footprint and computational efficiency}

Unlike the semi-implicit full PIC solver, which allocates auxiliary GMRES Krylov vectors, preconditioner blocks, and mass-matrix MultiFabs on each AMR level, the explicit full PIC solver operates directly on the core $(\mathbf{E}, \mathbf{B}, \mathbf{J})$ MultiFabs. This reduces peak memory usage per grid cell by approximately $60\%$ and avoids iterative linear solver convergence checks, making it well-suited for high-resolution relativistic and laser-plasma or electromagnetic shock regimes.

\section{Miscellaneous}
\subsection{Calculate total pressure tensor from sub groups}
Pressure tensor for each sub group is:
\begin{eqnarray}
    p^s_{m,n} = \frac{1}{V} \sum_{i=1} w_i^s (v_{i,m}^s-\bar{v}_m^s) (v_{i,n}^s-\bar{v}_n^s) 
\end{eqnarray}
where $V$ is the volume, and $\bar{v}_n^2$ and $\bar{v}_m^s$ are the average velocity of the sub group $s$. 

The total pressure tensor is:
\begin{eqnarray}
    p_{m,n} = \frac{1}{V} \sum_{s}\sum_{i=1} w_i^s (v_{i,m}^s-\bar{v}_m) (v_{i,n}^s-\bar{v}_n)\\
    \bar{v}_m = \frac{\sum_{s}\sum_{i=1} w_i^s v_{i,m}^s}{\sum_{s}\sum_{i=1} w_i^s} = \frac{\sum_s \rho_s \bar{v}_m}{\sum_{s}\rho_s} \\
    \bar{v}_n = \frac{\sum_{s}\sum_{i=1} w_i^s v_{i,n}^s}{\sum_{s}\sum_{i=1} w_i^s} = = \frac{\sum_s \rho_s \bar{v}_n}{\sum_{s}\rho_s}
\end{eqnarray}
where $\rho_s$ is the density of sub group $s$, $\bar{v}_n$ and $\bar{v}_m$ are the average velocity of all particles

\begin{eqnarray}
    p_{m,n} = \frac{1}{V} \sum_{s}\sum_{i=1} w_i^s [(v_{i,m}^s- \bar{v}_m^s) + (\bar{v}_m^s -\bar{v}_m)] [(v_{i,n}^s- \bar{v}_n^s) + (\bar{v}_n^s  - \bar{v}_n)]\\
    = \frac{1}{V} \sum_{s}\sum_{i=1} w_i^s[(v_{i,m}^s- \bar{v}_m^s)(v_{i,n}^s- \bar{v}_n^s)] + \\ \frac{1}{V} \sum_{s}\sum_{i=1} w_i^s[(v_{i,m}^s- \bar{v}_m^s)(\bar{v}_n^s - \bar{v}_n)] + \\
    \frac{1}{V} \sum_{s}\sum_{i=1} w_i^s[(\bar{v}_m^s -\bar{v}_m)(v_{i,n}^s- \bar{v}_n^s)] + \\ \frac{1}{V} \sum_{s}\sum_{i=1} w_i^s[(\bar{v}_m^s -\bar{v}_m)(\bar{v}_n^s - \bar{v}_n)]
\end{eqnarray}
The second and third terms are zeros. The first term is simply the sum of all the sub group tensors. The fourth term is $\sum_s \rho_s [(\bar{v}_m^s -\bar{v}_m)(\bar{v}_n^s - \bar{v}_n)]$. So, the final expression is:
\begin{eqnarray}
    p_{m,n} = \sum_s p^s_{m,n} + \sum_s \rho_s [(\bar{v}_m^s -\bar{v}_m)(\bar{v}_n^s - \bar{v}_n)]
\end{eqnarray}
\end{document}